% 03-modelo-sf-harris.tex — Chapter 3: El Modelo SF-Harris

\chapter{El Modelo SF-Harris}
\label{ch:modelo-sf-harris}

\section{Especificaci\'{o}n del modelo}

El modelo SF-Harris~\citep{anzarut2024} es un modelo de volatilidad estoc\'{a}stica con
recurrencia de Harris. La volatilidad latente $\tau_t$ sigue una cadena de Markov:

\begin{equation}
    \log(\tau_t) = \begin{cases}
        \log(\tau_{t-1}) & \text{con probabilidad } \pstay = e^{-\alpha} \quad \text{(estancia)} \\
        X_t \sim Q(\mu, \sigma) & \text{con probabilidad } 1 - \pstay \quad \text{(salto)}
    \end{cases}
\end{equation}

con emisi\'{o}n observacional:
\begin{equation}
    Y_t = \log(\RV_t) = \log(\tau_t) + \varepsilon_t, \quad \varepsilon_t \sim \N(0, \sigma_{\text{obs}}^2)
\end{equation}

donde $\RV_t$ es la varianza realizada intradiaria y $\sigma_{\text{obs}} \approx 1.04$
para datos de 15 minutos.

La distribuci\'{o}n invariante es $Q$, que Anzarut parametriza como $\GIG(\lambda, \kappa, \eta)$.
Cuando el proceso salta, extrae un nuevo nivel de volatilidad de $Q$;
cuando se queda, mantiene su valor anterior.

\section{La tesis de Anzarut}

Anzarut~\citep{anzarut2024} propone:
\begin{enumerate}
    \item Umbral $\varepsilon$ para detectar saltos: $|X_{t+1} - X_t| > \varepsilon$ indica un salto.
    \item Distribuci\'{o}n GIG param\'{e}trica para $Q$, estimada por m\'{a}xima verosimilitud.
    \item Muestreador de Gibbs para $(\alpha, \mu, \sigma)$ con priors conjugados.
    \item Ajuste de periodicidad intradiaria: $f(t) = \RV_{\text{15min}}(t) / \overline{\RV}$.
\end{enumerate}

La Tabla~3 de la tesis reporta AAD = 0.8\,pp para cobertura predictiva de $\log(\RV_t)$
en datos de IBM a frecuencia de 15 minutos, superando a GARCH(1,1) (3.4\,pp) y modelos Heston
($>$50\,pp).

\section{Replicaci\'{o}n: Resultados de la Tabla~3}

Replicamos la Tabla~3 con dos variantes de agregaci\'{o}n temporal:

\begin{table}[h]
\centering
\caption{Cobertura predictiva de $\log(\RV_t)$: SF-Harris vs.\ benchmarks.}
\label{tab:table3}
\begin{tabular}{lccccccc}
\toprule
Modelo & Barras & $p=0.25$ & $0.50$ & $0.75$ & $0.85$ & $0.90$ & $0.95$ & $\AAD$ \\
\midrule
SF-Harris Gibbs ($\varepsilon=0.1$) & D\'{o}lar & 25.1\% & 50.2\% & 75.0\% & 84.9\% & 90.0\% & 95.1\% & \textbf{0.2} \\
SF-Harris Gibbs ($\varepsilon=0.1$) & 15\,min & 25.2\% & 50.5\% & 75.3\% & 85.1\% & 90.2\% & 94.8\% & 0.3 \\
Anzarut (GIG+Gibbs) & 15\,min & 25\% & 51\% & 75\% & 84\% & 89\% & 93\% & 0.8 \\
GARCH(1,1) & 15\,min & 23.1\% & 46.6\% & 70.1\% & 80.6\% & 86.6\% & 92.7\% & 3.4 \\
Sim.\ Hist\'{o}rica & 15\,min & 25.0\% & 50.1\% & 75.2\% & 85.0\% & 90.0\% & 95.1\% & 0.4 \\
\bottomrule
\end{tabular}
\end{table}

Nuestro resultado de 0.2\,pp con barras de d\'{o}lar supera los 0.8\,pp originales por un
factor de $4\times$. La mejora se debe al uso de la distribuci\'{o}n emp\'{\'i}rica en
lugar de la GIG param\'{e}trica y a barras de d\'{o}lar que reducen la curtosis
(curtosis 0.09 vs 35.5 para barras de calendario).

\section{Barras de d\'{o}lar vs.\ barras de calendario}

\begin{table}[h]
\centering
\caption{Estad\'{\'i}sticas de $\log(\RV)$ por tipo de barra.}
\begin{tabular}{lcc}
\toprule
Estad\'{\'i}stica & 15\,min & D\'{o}lar \\
\midrule
Curtosis & 35.5 & 0.09 \\
Asimetr\'{\'i}a & $-$4.8 & 0.03 \\
Desv.\ est\'{a}ndar & 1.04 & 0.78 \\
AAD (Gibbs $\varepsilon=0.1$) & 0.3\,pp & 0.2\,pp \\
\bottomrule
\end{tabular}
\end{table}

Las barras de d\'{o}lar~\citep{lopezdeprado2018} muestrean a volumen constante en vez de
tiempo constante, lo que reduce la asimetr\'{\'i}a y curtosis de la distribuci\'{o}n de
$\log(\RV)$. Esto mejora la cobertura porque la distribuci\'{o}n emp\'{\'i}rica $Q$
captura mejor la forma marginal cuando los datos son menos extremos.